\FigureLayoutDeclare{img/2004/fujian_liberal/q16_fig1.png}{7.5cm}{19.999bp 15.199bp 37.198bp 7.6bp}
\FigureLayoutDeclare{img/2004/fujian_liberal/q21_fig1.png}{7.0cm}{135bp 34bp 42bp 64bp}

\chapter{2004年福建卷（文）}
\section{单选题}
\begin{problem}
设集合 \(U = \{1, 2, 3, 4, 5\}\)，\(A = \{1, 3, 5\}\)，\(B = \{2, 3, 5\}\)，则 \(\complement_U (A \cap B)\) 等于（\quad）

\choices
    {\(\{1,2,4\}\)}
    {\(\{4\}\)}
    {\(\{3,5\}\)}
    {\(\varnothing\)}
\end{problem}

\begin{answer}
A
\end{answer}

\begin{solution}
先求交集：\(A\cap B=\{3,5\}\). 补集是相对于全集 \(U\) 的；从全集中除去交集中的元素，得到
\(\complement_U(A\cap B)=\{1,2,4\}\)，故选 A.
\end{solution}

\begin{problem}
\(\tan 15^{\circ} + \cot 15^{\circ}\) 的值是（\quad）

\choices
    {\(2\)}
    {\(2 + \sqrt{3}\)}
    {\(4\)}
    {\(\frac{4\sqrt{3}}{3}\)}
\end{problem}

\begin{answer}
C
\end{answer}

\begin{solution}
因为 \(15^{\circ}\) 的正弦、余弦均非零，所以
\[
\tan15^{\circ}+\cot15^{\circ}
=\frac{\sin15^{\circ}}{\cos15^{\circ}}+\frac{\cos15^{\circ}}{\sin15^{\circ}}
=\frac{1}{\sin15^{\circ}\cos15^{\circ}}
\]
又 \(2\sin15^{\circ}\cos15^{\circ}=\sin30^{\circ}=\frac12\)，故
\(\sin15^{\circ}\cos15^{\circ}=\frac14\)，从而原式为 \(4\)，故选 C.
\end{solution}

\begin{problem}
命题 \(p\)：若 \(a\)，\(b \in \R\)，则 \(|a| + |b| > 1\) 是 \(|a + b| > 1\) 的充要条件. 命题 \(q\)：函数 \(y = \sqrt{|x - 1|} - 2\) 的定义域是 \((- \infty, -1] \cup [3, +\infty)\). 则（\quad）

\choices
    {“\(p\) 或 \(q\)”为假}
    {“\(p\) 且 \(q\)”为真}
    {\(p\) 真 \(q\) 假}
    {\(p\) 假 \(q\) 真}
\end{problem}

\begin{answer}
A
\end{answer}

\begin{solution}
取 \(a=\frac35\)，\(b=-\frac35\)，则 \(|a|+|b|=\frac65>1\)，但 \(|a+b|=0\)，所以条件 \(|a|+|b|>1\) 不能推出 \(|a+b|>1\)，命题 \(p\) 为假.

对任意 \(x\in\R\)，都有 \(|x-1|\geq0\)，所以 \(\sqrt{|x-1|}\) 有意义，函数 \(y=\sqrt{|x-1|}-2\) 的定义域是 \(\R\)，不是题中给出的 \((-\infty,-1]\cup[3,+\infty)\)，故命题 \(q\) 也为假. 因此“\(p\) 或 \(q\)”为假，故选 A.
\end{solution}

\begin{problem}
已知 \(F_{1}\)，\(F_{2}\) 是椭圆的两个焦点，过 \(F_{1}\) 且与椭圆长轴垂直的直线交椭圆于 \(A\)，\(B\) 两点，若 \(\triangle ABF_{2}\) 是正三角形，则这个椭圆的离心率是（\quad）

\choices
    {\(\frac{\sqrt{2}}{3}\)}
    {\(\frac{\sqrt{3}}{3}\)}
    {\(\frac{\sqrt{2}}{2}\)}
    {\(\frac{\sqrt{3}}{2}\)}
\end{problem}

\begin{answer}
\(\frac{\sqrt{3}}{3}\)
\end{answer}

\begin{solution}
设椭圆方程为 \(\frac{x^2}{a^2}+\frac{y^2}{b^2}=1\)，其中半焦距为 \(c\)，则 \(b^2=a^2-c^2\). 取 \(F_1=(-c,0)\)，\(F_2=(c,0)\)，过 \(F_1\) 且垂直于长轴的直线为 \(x=-c\). 代入椭圆方程，交点 \(A\)，\(B\) 的纵坐标满足
\(\frac{c^2}{a^2}+\frac{y^2}{b^2}=1\)，\(y^2=\frac{b^2(a^2-c^2)}{a^2}=\frac{b^4}{a^2}\).
因此 \(AB=\frac{2b^2}{a}\)，且由对称性知 \(F_1\) 是线段 \(AB\) 的中点.

三角形 \(ABF_2\) 是正三角形，故从顶点 \(F_2\) 到边 \(AB\) 的高为 \(F_1F_2=2c\)，同时等于 \(\frac{\sqrt3}{2}AB\). 于是
\[
2c=\frac{\sqrt3 b^2}{a}=\sqrt3\frac{a^2-c^2}{a}
\]
令离心率 \(e=\frac ca\)，得到 \(2e=\sqrt3(1-e^2)\)，即
\(\sqrt3e^2+2e-\sqrt3=0\). 解得 \(e=\frac{\sqrt3}{3}\) 或 \(e=-\sqrt3\)，因 \(e>0\)，所以 \(e=\frac{\sqrt3}{3}\)，故选 B.
\end{solution}

\begin{problem}
设 \(S_{n}\) 是等差数列 \(\{a_{n}\}\) 的前 \(n\) 项和，若 \(\frac{a_{5}}{a_{3}} = \frac{5}{9}\)，则 \(\frac{S_{9}}{S_{5}} =\)（\quad）

\choices
    {\(1\)}
    {\(-1\)}
    {\(2\)}
    {\(\frac{1}{2}\)}
\end{problem}

\begin{answer}
A
\end{answer}

\begin{solution}
设首项为 \(a_1\)，公差为 \(d\). 由题设分母不为零，且
\[
\frac{a_5}{a_3}=\frac{a_1+4d}{a_1+2d}=\frac59
\]
得 \(9a_1+36d=5a_1+10d\)，即 \(2a_1+13d=0\). 若 \(d=0\)，则题设比值为 \(1\)，矛盾，所以 \(d\neq0\).

等差数列前 \(n\) 项和为 \(S_n=\frac n2[2a_1+(n-1)d]=n\left(a_1+\frac{n-1}{2}d\right)\)，故
\[
\frac{S_9}{S_5}=\frac{9(a_1+4d)}{5(a_1+2d)}=\frac95\cdot\frac59=1
\]
因此选 A.
\end{solution}

\begin{problem}
已知 \(m\)，\(n\) 是不重合的直线，\(\alpha\)，\(\beta\) 是不重合的平面，有下列命题：

\begin{circlelist}
\item 若 \(m \subset \alpha\)，\(n \parallel \alpha\)，则 \(m \parallel n\)；
\item 若 \(m \parallel \alpha\)，\(m \parallel \beta\)，则 \(\alpha \parallel \beta\)；
\item 若 \(\alpha \cap \beta = n\)，\(m \parallel n\)，则 \(m \parallel \alpha\) 且 \(m \parallel \beta\)；
\item 若 \(m \perp \alpha\)，\(m \perp \beta\)，则 \(\alpha \parallel \beta\).
\end{circlelist}

其中真命题的个数是（\quad）

\choices
    {\(0\)}
    {\(1\)}
    {\(2\)}
    {\(3\)}
\end{problem}

\begin{answer}
B
\end{answer}

\begin{solution}
第一个命题不一定成立. 例如取 \(\alpha:z=0\)，令 \(m\) 为 \(x\) 轴，令 \(n\) 过点 \((0,0,1)\) 且方向向量为 \((1,1,0)\). 则 \(n\parallel\alpha\)，但 \(m\) 与 \(n\) 方向不相同且不相交，二者异面，所以 \(m\parallel n\) 不成立.

第二个命题也不一定成立. 取 \(\alpha:z=0\)，\(\beta:y=0\)，二者相交于 \(x\) 轴；令 \(m\) 为过 \((0,1,1)\)、方向向量为 \((1,0,0)\) 的直线. 此时 \(m\) 同时平行于 \(\alpha\) 和 \(\beta\)，但 \(\alpha\) 与 \(\beta\) 相交，故不能推出 \(\alpha\parallel\beta\).

第三个命题不一定成立. 仍取 \(\alpha:z=0\)，\(\beta:y=0\)，则 \(n=\alpha\cap\beta\) 为 \(x\) 轴. 令 \(m\) 为平面 \(\alpha\) 内与 \(n\) 平行且不同于 \(n\) 的直线，例如 \(y=1\)，\(z=0\). 此时 \(m\parallel n\)，但 \(m\subset\alpha\)，不称为 \(m\parallel\alpha\)，故结论不成立.

第四个命题成立. 同一直线垂直于两个平面，说明两个平面的法向方向都与该直线平行，因而两个平面的法向量平行. 两个平面又不重合，所以它们平行. 故只有一个真命题，选 B.
\end{solution}

\begin{problem}
已知函数 \(y = \log_{2} x\) 的反函数是 \(y = f^{-1}(x)\)，则函数 \(y = f^{-1}(1 - x)\) 的图象是（\quad）

\choices
    {\choicebitmap[width=0.15\paperwidth]{img/2004/fujian_liberal/q7_optA.png}}
    {\choicebitmap[width=0.15\paperwidth]{img/2004/fujian_liberal/q7_optB.png}}
    {\choicebitmap[width=0.15\paperwidth]{img/2004/fujian_liberal/q7_optC.png}}
    {\choicebitmap[width=0.15\paperwidth]{img/2004/fujian_liberal/q7_optD.png}}
\end{problem}

\begin{answer}
C
\end{answer}

\begin{solution}
函数 \(y=\log_2x\) 的反函数为 \(y=f^{-1}(x)=2^x\)，所以
\[
y=f^{-1}(1-x)=2^{1-x}=2\left(\frac12\right)^x
\]
底数 \(\frac12\in(0,1)\)，故该函数在 \(\R\) 上单调递减；又 \(x=0\) 时 \(y=2\)，所以图象过 \((0,2)\)，并以 \(x\) 轴为水平渐近线. 符合这些特征的图象是选项 C，故选 C.
\end{solution}

\begin{problem}
已知 \(\bs{a}\)，\(\bs{b}\) 是非零向量且满足 \((\bs{a} - 2\bs{b}) \perp \bs{a}\)，\((\bs{b} - 2\bs{a}) \perp \bs{b}\)，则 \(\bs{a}\) 与 \(\bs{b}\) 的夹角是（\quad）

\choices
    {\(\frac{\pi}{6}\)}
    {\(\frac{\pi}{3}\)}
    {\(\frac{2\pi}{3}\)}
    {\(\frac{5\pi}{6}\)}
\end{problem}

\begin{answer}
B
\end{answer}

\begin{solution}
由垂直条件和向量数量积的定义，得
\((\bs a-2\bs b)\cdot\bs a=0\)，\((\bs b-2\bs a)\cdot\bs b=0\).
即
\(|\bs a|^2=2\bs a\cdot\bs b\)，\(|\bs b|^2=2\bs a\cdot\bs b\).
于是 \(|\bs a|=|\bs b|\neq0\)，并且
\(\bs a\cdot\bs b=\frac12|\bs a|^2>0\). 若夹角为 \(\theta\)，则
\[
\cos\theta=\frac{\bs a\cdot\bs b}{|\bs a||\bs b|}
=\frac{\frac12|\bs a|^2}{|\bs a|^2}=\frac12
\]
由于 \(0\leq\theta\leq\pi\)，可得 \(\theta=\frac\pi3\)，故选 B.
\end{solution}

\begin{problem}
已知 \(\left(x - \frac{a}{x}\right)^{8}\) 展开式中常数项为 \(1120\)，其中实数 \(a\) 是常数，则展开式中各项系数的和是（\quad）

\choices
    {\(2^{8}\)}
    {\(3^{8}\)}
    {\(1\) 或 \(3^{8}\)}
    {\(1\) 或 \(2^{8}\)}
\end{problem}

\begin{answer}
C
\end{answer}

\begin{solution}
展开式通项为
\(\mathrm C_8^k x^{8-k}\left(-\frac ax\right)^k =\mathrm C_8^k(-a)^k x^{8-2k}\)，\(k=0\)，\(1\)，\(\ldots\)，\(8\).
常数项要求 \(8-2k=0\)，故只能取 \(k=4\). 于是
\(\mathrm C_8^4a^4=1120\)，即 \(70a^4=1120\)，所以 \(a^4=16\)，实数 \(a\) 只能为 \(2\) 或 \(-2\).

多项式各项系数的和等于令 \(x=1\) 时的值，因此系数和为 \((1-a)^8\). 当 \(a=2\) 时为 \(1\)，当 \(a=-2\) 时为 \(3^8\)，故选 C.
\end{solution}

\begin{problem}
如图，\(A\)，\(B\)，\(C\) 是表面积为 \(48\pi\) 的球面上三点，\(AB = 2\)，\(BC = 4\)，\(\angle ABC = 60^{\circ}\)，\(O\) 为球心，则直线 \(OA\) 与截面 \(ABC\) 所成的角是（\quad）

\bitmapfigure[max width=0.42\linewidth]{img/2004/fujian_liberal/q10_fig1.png}

\choices
    {\(\arcsin \frac{\sqrt{3}}{6}\)}
    {\(\arccos \frac{\sqrt{3}}{6}\)}
    {\(\arcsin \frac{\sqrt{3}}{3}\)}
    {\(\arccos \frac{\sqrt{3}}{3}\)}
\end{problem}

\begin{answer}
D
\end{answer}

\begin{solution}
设球半径为 \(R\). 由球的表面积得
\(4\pi R^2=48\pi\)，所以 \(R=2\sqrt3\).

在 \(\triangle ABC\) 中，由余弦定理
\[
AC^2=AB^2+BC^2-2\cdot AB\cdot BC\cos60^{\circ}
=2^2+4^2-2\cdot2\cdot4\cdot\frac12=12
\]
所以 \(AC=2\sqrt3\). 设三角形 \(ABC\) 的外接圆半径为 \(r\)，由正弦定理
\[
r=\frac{AC}{2\sin\angle ABC}
=\frac{2\sqrt3}{2\sin60^{\circ}}=2
\]
球心 \(O\) 到平面 \(ABC\) 的垂足 \(H\) 是该外接圆圆心，因此 \(HA=r=2\).

设直线 \(OA\) 与平面 \(ABC\) 所成的角为 \(\theta\). 在直角三角形 \(OHA\) 中，\(OA=R=2\sqrt3\)，且 \(HA\) 是 \(OA\) 在平面内的投影，故
\[
\cos\theta=\frac{HA}{OA}=\frac{2}{2\sqrt3}=\frac{\sqrt3}{3}
\]
所以 \(\theta=\arccos\frac{\sqrt3}{3}\)，故选 D.
\end{solution}

\begin{problem}
定义在 \(\R\) 上的偶函数 \(f(x)\) 满足 \(f(x) = f(x + 2)\)，当 \(x \in [3,4]\) 时，\(f(x) = x - 2\)，则（\quad）

\choices
    {\(f\left(\sin \frac{1}{2}\right) < f\left(\cos \frac{1}{2}\right)\)}
    {\(f\left(\sin \frac{\pi}{3}\right) > f\left(\cos \frac{\pi}{3}\right)\)}
    {\(f(\sin 1) < f(\cos 1)\)}
    {\(f\left(\sin \frac{3}{2}\right) > f\left(\cos \frac{3}{2}\right)\)}
\end{problem}

\begin{answer}
C
\end{answer}

\begin{solution}
令 \(|t|\leq1\). 由偶性，\(f(t)=f(|t|)=f(-|t|)\)；再由周期为 \(2\)，有
\(f(-|t|)=f(2-|t|)=f(4-|t|)\)，其中 \(4-|t|\in[3,4]\). 利用已知函数值，得到
\[
f(t)=f(4-|t|)=(4-|t|)-2=2-|t|
\]

因为 \(0<\frac12<\frac\pi4\)，所以 \(\sin\frac12<\cos\frac12\)，从而 \(f(\sin\frac12)>f(\cos\frac12)\)，A 错. 又 \(\frac\pi3>\frac\pi4\)，故 \(\sin\frac\pi3>\cos\frac\pi3\)，所以 B 中不等号方向相反，B 错. 由于 \(\frac\pi4<1<\frac\pi2\)，有 \(\sin1>\cos1\)，从而
\(f(\sin1)=2-\sin1<2-\cos1=f(\cos1)\)，C 对. 最后 \(\frac\pi4<\frac32<\frac\pi2\)，有 \(\sin\frac32>\cos\frac32\)，故 D 错. 因此选 C.
\end{solution}

\begin{problem}
如图，\(B\) 地在 \(A\) 地的正东方向 \(4\mathrm{~km}\) 处，\(C\) 地在 \(B\) 地的北偏东 \(30^{\circ}\) 方向 \(2\mathrm{~km}\) 处，河流的某岸 \(PQ\)（曲线）上任意一点到 \(A\) 的距离比到 \(B\) 的距离远 \(2\mathrm{~km}\). 现要在曲线 \(PQ\) 上选一处 \(M\) 建一座码头，向 \(B\)，\(C\) 两地转运货物. 经测算，从 \(M\) 到 \(B\)，\(C\) 两地修建公路的费用分别是 \(a\text{ 万元}/\mathrm{km}\)，\(2a\text{ 万元}/\mathrm{km}\)，那么修建这两条公路的总费用最低是（\quad）

\bitmapfigure[max width=0.34\linewidth]{img/2004/fujian_liberal/q12_fig1.png}

\choices
    {\((\sqrt{7} + 1)a\text{ 万元}\)}
    {\((2\sqrt{7} - 2)a\text{ 万元}\)}
    {\(2\sqrt{7}a\text{ 万元}\)}
    {\((\sqrt{7} - 1)a\text{ 万元}\)}
\end{problem}

\begin{answer}
按题面文字计算为 \(5a\)，但该值不在所列四个选项中.
\end{answer}

\begin{solution}
取 \(A=(0,0)\)，\(B=(4,0)\)，令正东方向为 \(x\) 轴正向、正北方向为 \(y\) 轴正向. 因为 \(BC=2\)，且 \(BC\) 为北偏东 \(30^{\circ}\) 方向，所以
\[
C=(4+2\sin30^{\circ},2\cos30^{\circ})=(5,\sqrt3)
\]

河岸上的点 \(M\) 满足 \(MA-MB=2\). 以 \(AB\) 的中点为原点，令 \(r=x-2\)，则由双曲线定义，\(M\) 在以 \(A\)，\(B\) 为焦点、实半轴长为 \(1\) 的右支上，故
\(\frac{r^2}{1^2}-\frac{y^2}{3}=1\)，\(r\geq1\)，\(y^2=3(r^2-1)\).
这里 \(x=r+2\). 于是
\[
MB^2=(x-4)^2+y^2=(r-2)^2+3(r^2-1)=(2r-1)^2
\]
因 \(r\geq1\)，所以 \(MB=2r-1\).

公路总费用除以正数 \(a\) 后为
\(MB+2MC\). 又
\[
MC^2=(x-5)^2+(y-\sqrt3)^2=(r-3)^2+(y-\sqrt3)^2
\]
当 \(1\leq r<3\) 时，距离不小于横坐标差的绝对值，故 \(MC\geq3-r\)，从而
\[
MB+2MC\geq(2r-1)+2(3-r)=5
\]
当 \(r\geq3\) 时，\(MC\geq0\)，故
\(MB+2MC\geq MB=2r-1\geq5\)，因此总费用至少为 \(5a\).

当 \(r=\sqrt2\)，\(y=\sqrt3\) 时，点 \(M=(2+\sqrt2,\sqrt3)\) 在上述双曲线上，且 \(M\)，\(C\) 同高，所以 \(MC=3-r\). 此时
\[
MB+2MC=(2\sqrt2-1)+2(3-\sqrt2)=5
\]
达到下界. 因此按原题文字，最低总费用为 \(5a\)，但原 \(\text{PDF}\) 所列四个选项均不等于 \(5a\).
\end{solution}

\section{填空题}
\begin{problem}
直线 \(x + 2y = 0\) 被曲线 \(x^{2} + y^{2} - 6x - 2y - 15 = 0\) 所截得的弦长等于\fillinblank{}.
\end{problem}

\begin{answer}
\(4\sqrt{5}\)
\end{answer}

\begin{solution}
将曲线方程配方，得
\[
x^2-6x+y^2-2y-15=0
\quad\Longleftrightarrow\quad
(x-3)^2+(y-1)^2=25
\]
所以圆心为 \(O_1=(3,1)\)，半径为 \(R=5\). 圆心到直线 \(x+2y=0\) 的距离为
\[
d=\frac{|3+2\cdot1|}{\sqrt{1^2+2^2}}=\sqrt5<5
\]
因此该直线确实与圆相交，弦长为
\[
2\sqrt{R^2-d^2}=2\sqrt{25-5}=2\sqrt{20}=4\sqrt5
\]
\end{solution}

\begin{problem}
设函数 \(f(x) = \begin{cases}
\frac{1}{2}x - 1, & (x \geqslant 0),\\
\frac{1}{x}, & (x < 0).
\end{cases}\) 若 \(f(a) > a\)，则实数 \(a\) 的取值范围是\fillinblank{}.
\end{problem}

\begin{answer}
\((-\infty,-1)\)
\end{answer}

\begin{solution}
按分段点 \(0\) 分类讨论.

当 \(a\geq0\) 时，\(f(a)=\frac12a-1\)，所以
\[
f(a)>a\quad\Longleftrightarrow\quad \frac12a-1>a
\quad\Longleftrightarrow\quad a<-2
\]
这与 \(a\geq0\) 矛盾，故这一段没有解.

当 \(a<0\) 时，\(f(a)=\frac1a\). 由于乘以负数 \(a\) 时不等号方向改变，
\[
\frac1a>a\quad\Longleftrightarrow\quad 1<a^2
\]
结合 \(a<0\)，得到 \(a<-1\). 端点 \(a=-1\) 只能得到等号，不满足严格不等式. 故实数 \(a\) 的取值范围为 \((-\infty,-1)\).
\end{solution}

\begin{problem}
一个总体中有 \(100\) 个个体，随机编号 \(0\)，\(1\)，\(2\)，\(\dots\)，\(99\)，依编号顺序平均分成 \(10\) 个小组，组号依次为 \(1\)，\(2\)，\(3\)，\(\dots\)，\(10\)，现用系统抽样方法抽取一个容量为 \(10\) 的样本. 规定如果在第 \(1\) 组随机抽取的号码为 \(m\)，那么在第 \(k\) 组中抽取的号码个位数字与 \(m + k\) 的个位数字相同. 若 \(m = 6\)，则在第 \(7\) 组中抽取的号码是\fillinblank{}.
\end{problem}

\begin{answer}
63
\end{answer}

\begin{solution}
每组有 \(10\) 个连续编号，第 \(7\) 组的编号为 \(60\)，\(61\)，\(\ldots\)，\(69\). 题设规定第 \(k\) 组抽取号码的个位数字与 \(m+k\) 的个位数字相同. 取 \(k=7\)，且 \(m=6\)，则
\(m+k=6+7=13\)，其个位数字为 \(3\). 因此第 \(7\) 组抽取的号码是 \(63\).
\end{solution}

\begin{problem}
如图，将边长为 \(1\) 的正六边形铁皮的六个角各切去一个全等的四边形，再沿虚线折起，做成一个无盖的正六棱柱容器. 当这个正六棱柱容器的底面边长为\fillinblank{} 时，其容积最大.

\bitmapfigure[max width=0.32\linewidth]{img/2004/fujian_liberal/q16_fig1.png}
\end{problem}

\begin{answer}
\(\frac{2}{3}\)
\end{answer}

\begin{solution}
设折成的正六棱柱底面边长为 \(x\). 由图形实际存在，\(0<x<1\). 原正六边形与折成的底面六边形对应边平行；正六边形边长每减少 \(1-x\)，相邻对应平行边的距离为 \(\frac{\sqrt3}{2}(1-x)\)，故棱柱高为
\[
h=\frac{\sqrt3}{2}(1-x)
\]
边长为 \(x\) 的正六边形面积为六个正三角形面积之和，即
\[
S=6\cdot\frac{\sqrt3}{4}x^2=\frac{3\sqrt3}{2}x^2
\]
所以容积函数为
\[
V(x)=Sh=\frac{3\sqrt3}{2}x^2\cdot\frac{\sqrt3}{2}(1-x)
=\frac94x^2(1-x)
\]
求导得
\[
V'(x)=\frac94(2x-3x^2)=\frac94x(2-3x)
\]
在 \(0<x<1\) 内，\(V'(x)>0\) 当 \(0<x<\frac23\)，\(V'(x)<0\) 当 \(\frac23<x<1\). 因此容积先增后减，在 \(x=\frac23\) 时取得最大值. 填 \(\frac23\).
\end{solution}

\section{解答题}
\begin{problem}
设函数 \(f(x) = \bs{a} \cdot \bs{b}\)，其中向量 \(\bs{a} = (2\cos x, 1)\)，\(\bs{b} = (\cos x, \sqrt{3}\sin 2x)\)，\(x \in \R\).

\begin{enumerate}
\item 若 \(f(x) = 1 - \sqrt{3}\) 且 \(x \in \left[-\frac{\pi}{3}, \frac{\pi}{3}\right]\)，求 \(x\)；
\item 若函数 \(y = 2\sin 2x\) 的图象按向量 \(\bs{c} = (m, n)\)（\(|m| < \frac{\pi}{2}\)）平移后得到函数 \(y = f(x)\) 的图象，求实数 \(m\)，\(n\) 的值.
\end{enumerate}
\end{problem}

\begin{solution}
由点积公式，
\[
f(x)=2\cos^2x+\sqrt3\sin2x
=1+\cos2x+\sqrt3\sin2x
=1+2\sin\left(2x+\frac{\pi}{6}\right)
\]

\begin{enumerate}
\item 当 \(f(x)=1-\sqrt3\) 时，
\[
\sin\left(2x+\frac{\pi}{6}\right)=-\frac{\sqrt3}{2}
\]
由 \(x\in\left[-\frac{\pi}{3},\frac{\pi}{3}\right]\)，可知
\[
2x+\frac{\pi}{6}\in\left[-\frac{\pi}{2},\frac{5\pi}{6}\right]
\]
在此区间内，正弦值为 \(-\frac{\sqrt3}{2}\) 的唯一角是 \(-\frac{\pi}{3}\)，故
\(2x+\frac{\pi}{6}=-\frac{\pi}{3}\)，\(x=-\frac{\pi}{4}\).
\item 将函数 \(y=2\sin2x\) 按向量 \((m,n)\) 平移后，横坐标为 \(x\) 的点来自原图象横坐标 \(x-m\) 的点，所以新图象方程为
\[
y=2\sin\bigl(2(x-m)\bigr)+n
=2\sin(2x-2m)+n
\]
要与 \(f(x)=1+2\sin(2x+\frac{\pi}{6})\) 恒相同，比较中线得 \(n=1\)，比较相位得
\[
-2m\equiv\frac{\pi}{6}\pmod{2\pi}
\]
因此 \(m=-\frac{\pi}{12}-k\pi\)，其中 \(k\in\Z\). 条件 \(|m|<\frac{\pi}{2}\) 只允许 \(k=0\)，故 \(m=-\frac{\pi}{12}\)，\(n=1\).
\end{enumerate}
\end{solution}

\begin{problem}
甲，乙两人参加一次英语口语考试，已知在备选的 \(10\) 道试题中，甲能答对其中的 \(6\) 题，乙能答对其中的 \(8\) 题. 规定每次考试都从备选题中随机抽出 \(3\) 题进行测试，至少答对 \(2\) 题才算合格.

\begin{enumerate}
\item 求甲，乙两人考试合格的概率；
\item 求甲，乙两人至少有一人考试合格的概率.
\end{enumerate}
\end{problem}

\begin{solution}
\begin{enumerate}
\item 每人的试题都从 \(10\) 道题中抽取 \(3\) 道，等可能的抽法总数为 \(\mathrm C_{10}^3=120\). 甲答对至少 \(2\) 题包括答对恰好 \(2\) 题和恰好 \(3\) 题两种情况，因此
\[
P(\text{甲合格})
=\frac{\mathrm C_6^2\mathrm C_4^1+\mathrm C_6^3}{\mathrm C_{10}^3}
=\frac{15\cdot4+20}{120}=\frac23
\]
乙能答对 \(8\) 道、答错 \(2\) 道，同理
\[
P(\text{乙合格})
=\frac{\mathrm C_8^2\mathrm C_2^1+\mathrm C_8^3}{\mathrm C_{10}^3}
=\frac{28\cdot2+56}{120}=\frac{14}{15}
\]
\item 两人的试题分别随机抽取，所以两人的合格事件相互独立. 两人都不合格的概率为
\[
\left(1-\frac23\right)\left(1-\frac{14}{15}\right)
=\frac13\cdot\frac1{15}=\frac1{45}
\]
故至少一人合格的概率为
\[
1-\frac1{45}=\frac{44}{45}
\]
\end{enumerate}
\end{solution}

\begin{problem}
在三棱锥 \(S - ABC\) 中，\(\triangle ABC\) 是边长为 \(4\) 的正三角形，平面 \(SAC \perp\) 平面 \(ABC\)，\(SA = SC = 2\sqrt{2}\)，\(M\) 为 \(AB\) 的中点.

\begin{enumerate}
\item 证明：\(AC \perp SB\)；
\item 求二面角 \(S - CM - B\) 的大小；
\item 求点 \(B\) 到平面 \(SCM\) 的距离.
\end{enumerate}

\bitmapfigure[max width=0.34\linewidth]{img/2004/fujian_liberal/q19_fig1.png}
\end{problem}

\begin{solution}
建立空间直角坐标系，取
\(A=(0,0,0)\)，\(C=(4,0,0)\)，\(B=(2,2\sqrt3,0)\).
于是平面 \(ABC\) 是 \(z=0\). 平面 \(SAC\) 与平面 \(ABC\) 垂直，且包含 \(AC\)，可取为 \(y=0\). 由 \(SA=SC\)，点 \(S\) 在 \(AC\) 的中垂面内；结合 \(S\in\{y=0\}\)，可设 \(S=(2,0,h)\). 由 \(SA=2\sqrt2\)，得
\((2-0)^2+h^2=8\)，\(h^2=4\).
取图示方向 \(h=2\)，所以 \(S=(2,0,2)\). 又 \(M\) 是 \(AB\) 的中点，故 \(M=(1,\sqrt3,0)\).

\begin{enumerate}
\item \(\overrightarrow{AC}=(4,0,0)\)，\(\overrightarrow{SB}=B-S=(0,2\sqrt3,-2)\). 因此
\[
\overrightarrow{AC}\cdot\overrightarrow{SB}=4\cdot0+0\cdot2\sqrt3+0\cdot(-2)=0
\]
所以 \(AC\perp SB\).
\item 取 \(C\) 为棱 \(CM\) 上的点，把 \(\overrightarrow{CS}\)，\(\overrightarrow{CB}\) 分别正投影到过 \(C\) 且垂直于 \(CM\) 的平面内. 先有
\(\overrightarrow{CM}=(-3,\sqrt3,0)\)，\(\overrightarrow{CS}=(-2,0,2)\)，\(\overrightarrow{CB}=(-2,2\sqrt3,0)\)，\(|\overrightarrow{CM}|^2=12\).
计算得 \(\overrightarrow{CS}\cdot\overrightarrow{CM}=6\)，\(\overrightarrow{CB}\cdot\overrightarrow{CM}=12\)，故两条投影向量分别为
\[
\overrightarrow{u}=\overrightarrow{CS}-\frac6{12}\overrightarrow{CM}
=\left(-\frac12,-\frac{\sqrt3}{2},2\right)
\]
\[
\overrightarrow{v}=\overrightarrow{CB}-\frac{12}{12}\overrightarrow{CM}
=(1,\sqrt3,0)
\]
它们分别位于半平面 \(SCM\) 和 \(BCM\) 内，且都垂直于棱 \(CM\)，所以 \(\angle(\overrightarrow u,\overrightarrow v)\) 是所求二面角的平面角. 又
\(\overrightarrow u\cdot\overrightarrow v=-2\)，\(|\overrightarrow u|=\sqrt5\)，\(|\overrightarrow v|=2\).
设二面角为 \(\theta\)，则
\[
\cos\theta=\frac{-2}{2\sqrt5}=-\frac1{\sqrt5}
\]
因此 \(\theta=\arccos\left(-\frac1{\sqrt5}\right)=\pi-\arctan 2\).
\item 由
\[
\overrightarrow{CM}\times\overrightarrow{CS}=(2\sqrt3,6,2\sqrt3)
\]
可取平面 \(SCM\) 的法向量为 \(\bs n=(\sqrt3,3,\sqrt3)\). 平面过点 \(C=(4,0,0)\)，故其方程为
\[
\sqrt3(x-4)+3y+\sqrt3z=0
\]
即 \(x+\sqrt3y+z-4=0\). 点 \(B=(2,2\sqrt3,0)\) 到此平面的距离为
\[
\frac{|2+\sqrt3\cdot2\sqrt3+0-4|}{\sqrt{1+3+1}}
=\frac4{\sqrt5}=\frac{4\sqrt5}{5}
\]
\end{enumerate}
\end{solution}

\begin{problem}
某企业 \(2003\) 年的纯利润为 \(500\text{ 万元}\)，因设备老化等原因，企业的生产能力将逐年下降. 若不能进行技术改造，预测从今年起每年比上一年纯利润减少 \(20\text{ 万元}\)，今年初该企业一次性投入资金 \(600\text{ 万元}\) 进行技术改造，预测在未扣除技术改造资金的情况下，第 \(n\) 年（今年为第一年）的利润为 \(500\left(1 + \frac{1}{2^{n}}\right)\text{ 万元}\)（\(n\) 为正整数）.

\begin{enumerate}
\item 设从今年起的前 \(n\) 年，若该企业不进行技术改造的累计纯利润为 \(A_n\text{ 万元}\)，进行技术改造后的累计纯利润为 \(B_n\text{ 万元}\)（须扣除技术改造资金），求 \(A_n\)，\(B_n\) 的表达式；
\item 依上述预测，从今年起该企业至少经过多少年，进行技术改造后的累计纯利润超过不进行技术改造的累计纯利润？
\end{enumerate}
\end{problem}

\begin{solution}
\begin{enumerate}
\item 若不进行技术改造，第 \(k\) 年的纯利润比第 1 年少 \(20(k-1)\) 万元，即为 \(500-20(k-1)\) 万元. 因此
\[
A_n=\sum_{k=1}^{n}[500-20(k-1)]
=500n-20\cdot\frac{n(n-1)}2
=510n-10n^2
\]
进行技术改造后，第 \(k\) 年的利润为 \(500(1+2^{-k})\) 万元，且今年年初一次性扣除 \(600\) 万元，所以
\[
B_n=\sum_{k=1}^{n}500\left(1+\frac1{2^k}\right)-600
=500n+500\sum_{k=1}^{n}\frac1{2^k}-600
\]
又
\[
B_n=500n+500(1-2^{-n})-600
=500n-100-\frac{500}{2^n}
\]
\item 令 \(D_n=B_n-A_n\)，则
\[
D_n=10n^2-10n-100-\frac{500}{2^n}
\]
对正整数 \(n\)，
\[
D_{n+1}-D_n=20n+\frac{250}{2^n}>0
\]
所以 \(\{D_n\}\) 严格递增. 计算相邻两项得
\[
D_4=10\cdot4^2-10\cdot4-100-\frac{500}{2^4}<0
\]
\[
D_5=10\cdot5^2-10\cdot5-100-\frac{500}{2^5}>0
\]
因 \(D_n\) 递增，最小的满足 \(D_n>0\) 的正整数是 \(n=5\). 所以至少经过 \(5\) 年，技术改造后的累计纯利润才超过不进行技术改造的累计纯利润.
\end{enumerate}
\end{solution}

\begin{problem}
如图，\(P\) 是抛物线 \(C: y = \frac{1}{2}x^2\) 上一点，直线 \(l\) 过点 \(P\) 且与抛物线 \(C\) 在点 \(P\) 的切线垂直，\(l\) 与抛物线 \(C\) 相交于另一点 \(Q\).

\begin{enumerate}
\item 当点 \(P\) 的横坐标为 \(2\) 时，求直线 \(l\) 的方程；
\item 当点 \(P\) 在抛物线 \(C\) 上移动时，求线段 \(PQ\) 中点 \(M\) 的轨迹方程，并求点 \(M\) 到 \(x\) 轴的最短距离.
\end{enumerate}

\bitmapfigure[max width=0.42\linewidth]{img/2004/fujian_liberal/q21_fig1.png}
\end{problem}

\begin{solution}
设 \(P=(p,\frac12p^2)\). 当 \(p=0\) 时，抛物线在原点处的切线为 \(x\) 轴，过原点的垂线为 \(x=0\)，它与抛物线没有另一个交点，因此题设有另一交点 \(Q\) 时必有 \(p\neq0\).

由 \(y=\frac12x^2\) 得 \(y'=x\)，所以在 \(P\) 处的切线斜率为 \(p\)，直线 \(l\) 的斜率为 \(-\frac1p\)，故
\[
y-\frac12p^2=-\frac1p(x-p)
\]

\begin{enumerate}
\item 当 \(p=2\) 时，\(P=(2,2)\)，代入直线方程得
\[
y-2=-\frac12(x-2)
\]
整理为 \(x+2y-6=0\).
\item 设 \(Q=(q,\frac12q^2)\). 将直线方程与抛物线方程联立，得
\(\frac12x^2=-\frac1p x+1+\frac12p^2\)，\(px^2+2x-p(p^2+2)=0\).
其中一个根是 \(p\)，由韦达定理，另一个根满足
\(p+q=-\frac2p\)，\(q=-p-\frac2p\).
设 \(M=(X,Y)\)，则
\[
X=\frac{p+q}{2}=-\frac1p
\]
且
\[
Y=\frac{\frac12p^2+\frac12q^2}{2}=\frac{p^2+q^2}{4}
=\frac12p^2+1+\frac1{p^2}
\]
由 \(X=-\frac1p\neq0\)，有 \(p^2=\frac1{X^2}\)，\(\frac1{p^2}=X^2\)，所以中点轨迹为
\(Y=X^2+1+\frac1{2X^2}\)，\(X\neq0\).
令 \(u=X^2>0\)，则
\[
Y=1+u+\frac1{2u}\geq1+2\sqrt{\frac{u}{2u}}=1+\sqrt2
\]
其中等号当且仅当 \(u=\frac1{2u}\)，即 \(X^2=\frac1{\sqrt2}\). 由于轨迹上 \(Y>0\)，点 \(M\) 到 \(x\) 轴的距离就是 \(Y\)，所以最短距离为 \(1+\sqrt2\).
\end{enumerate}
\end{solution}

\begin{problem}
已知 \(f(x) = 4x + ax^2 - \frac{2}{3}x^3\)（\(x \in \R\)）在区间 \([-1,1]\) 上是增函数.

\begin{enumerate}
\item 求实数 \(a\) 的值组成的集合 \(A\)；
\item 设关于 \(x\) 的方程 \(f(x) = 2x + \frac{1}{3}x^3\) 的两个非零实根为 \(x_1\)，\(x_2\). 试问：是否存在实数 \(m\)，使得不等式 \(m^2 + tm + 1 \geqslant |x_1 - x_2|\) 对任意 \(a \in A\) 及 \(t \in [-1,1]\) 恒成立？若存在，求 \(m\) 的取值范围；若不存在，请说明理由.
\end{enumerate}
\end{problem}

\begin{solution}
\begin{enumerate}
\item 先求导：
\[
f'(x)=4+2ax-2x^2
\]
因为 \(f\) 在闭区间上可导，\(f\) 为增函数等价于 \(f'(x)\geq0\) 对所有 \(x\in[-1,1]\) 成立. 对 \(x\in(0,1]\)，不等式化为
\[
4+2ax-2x^2\geq0
\quad\Longleftrightarrow\quad
a\geq x-\frac2x
\]
函数 \(x-\frac2x\) 在 \((0,1]\) 上的导数为 \(1+\frac2{x^2}>0\)，故其最大值为在 \(x=1\) 处的 \(-1\)，从而需要 \(a\geq-1\).

对 \(x\in[-1,0)\)，除以负数 \(x\) 后不等号方向改变，得
\[
a\leq x-\frac2x
\]
同一函数在 \([-1,0)\) 上递增，其最小值为在 \(x=-1\) 处的 \(1\)，从而需要 \(a\leq1\). 在 \(x=0\) 处 \(f'(0)=4>0\)，不产生额外限制. 因此
\[
A=[-1,1]
\]
\item 原方程等价于
\[
4x+ax^2-\frac23x^3=2x+\frac13x^3
\quad\Longleftrightarrow\quad
x(2+ax-x^2)=0
\]
除去非零根 \(x=0\)，两个根满足
\[
x^2-ax-2=0
\]
其判别式为 \(a^2+8>0\)，所以两个非零实根为
\(x_1=\frac{a+\sqrt{a^2+8}}2\)，\(x_2=\frac{a-\sqrt{a^2+8}}2\)
从而
\[
|x_1-x_2|=\sqrt{a^2+8}
\]
当 \(a\in[-1,1]\) 时，\(a^2\leq1\)，故 \(|x_1-x_2|\leq3\)，且在 \(a=\pm1\) 时取得最大值 \(3\).

固定 \(m\)，当 \(t\in[-1,1]\) 变化时，线性式 \(tm\) 的最小值为 \(-|m|\)，所以
\[
\min_{-1\leq t\leq1}(m^2+tm+1)=m^2+1-|m|
\]
要使不等式对所有 \(a\in A\) 和 \(t\in[-1,1]\) 恒成立，充要条件是
\[
m^2+1-|m|\geq3
\]
令 \(u=|m|\geq0\)，则
\[
u^2-u-2=(u-2)(u+1)\geq0
\]
由于 \(u+1>0\)，得 \(u\geq2\). 因此这样的实数存在，且
\[
m\in(-\infty,-2]\cup[2,+\infty)
\]
\end{enumerate}
\end{solution}
